\begin{topic}[id=canopen,
             difficulty={Mittel},
             type={Protokollvergleich},
             prerequisites={Grundlagen eingebetteter Systeme; Basiswissen Feldbusse hilfreich},
             practice={optional},
             owner={Dozententeam},
             updated={2026-04-09},
             tags={ CANopen, CiA, PDO, SDO, EDS/DCF, NMT, Diagnose }]{CANopen – Protokoll \& Praxisvergleich}

\TopicDescription{CANopen baut auf dem robusten Feldbus CAN auf und bietet Profile, Objektverzeichnisse und Mechanismen für Konfiguration und zyklische Daten (PDO/SDO). Du lernst, wie EDS/DCF-Dateien genutzt werden, was NMT-States bedeuten und wie Diagnose/Fehlerhandling funktioniert. Wir ordnen CANopen im Vergleich zu modernen Ethernet-basierten Protokollen ein (Profinet, OPC UA over TSN) und diskutieren Migrationspfade in Brownfield-Anlagen. Ziel ist ein realistisches Verständnis, wann CANopen weiterhin sinnvoll ist – und wie man Mischumgebungen betreibt.}

\TopicLeitfragen{\item Welche Rolle spielen PDO vs. SDO im Design?
\item Wie gelingt Diagnose/Fehlerhandling pragmatisch?
\item Wann lohnt eine Migration auf Ethernet-Protokolle?
}
\TopicScope{\item Keine Stecker-/Leitungstechnik.
\item Keine tiefe TSN-Konfiguration.
}
\TopicPractice{Optionale Praxisidee: Kleines Objektverzeichnis lesen/schreiben und PDO-Mapping zeigen.}

\TopicKeywords{CANopen, CiA, PDO, SDO, EDS/DCF, NMT, Diagnose}

\nocite{ciaCanopen,vectorCanopen}
\end{topic}
